% Q2.tex
\section{Question 2}
\begin{enumerate}[label=(\alph*)]
    \item 
    \begin{enumerate}[label=(\roman*)]
        \item 
        Let $ \mathcal{R} $ be the simply connected region $\{ x: |z| < 3/2  \}$,  \\
        $f(z) = \dfrac{ e^{-2z} }{ z+2i }$. \\
        Then $f$ is analytic on $ \mathcal{R} $ and $ C $ is a simple-connected contour in $\mathcal{R} $.
        \vs \\
        Thus by CAUCHY'S THEOREM,
        \[ \int_C = \frac{ e^{-2z} }{ z+2i } = 0. \]
        \vvs \\
        \item 
        Let $ \mathcal{R} $ and $f$ be as above, meeting the same conditions.
        \vs \\
        Then, by CAUCHY'S INTEGRAL FORMULA,
        \begin{align*}
            \int_C \frac{ e^{-2z} }{ z(z+2i) } &= 2\pi i f(0) \\
            &= 2\pi i \left( \frac{1}{0 + 21} \right) = \pi.
        \end{align*}
        \vvs \\
        \item 
        With $ \mathcal{R} $ and $f$ as above, meeting the same conditions,
        \vs \\
        by CAUCHY'S FIRST DERIVATIVE FORMULA, as
        \begin{align*}
            f'(z) &= \frac{ (z+2i)(-2e^{-2z}) - e^{-2z}(1) }{ (z+2i)^2 }
            \intertext{by the QUOTIENT RULE for differentiation,}
            &= - \frac{ e^{-2z}( 2z+4i+1 ) }{(z+2i)^2} \\
            \int_C \frac{ e^{-2z} }{ z2(z+2i) } &= 2\pi i f'(0) \\
            &= 2\pi i \left( -\frac{ 1+4i }{ (2i)^2 } \right) \\
            &= 2\pi i \left( \frac{ 1+4i }{ 4 } \right) \\
            &= -2\pi + \frac{1}{2}\pi i = \pi \left( -2 + \frac{1}{2}i \right).
        \end{align*}
    \end{enumerate}
    \leavevmode \vvvs \\
    \item 
    $C$ is the unit circle and $\Gamma$ is the circle of radius 1/2 centred at the origin but traversed in the opposite direction. It also nominally has different start and end points, but this is irrelevant as it is a CLOSED CONTOUR.
    \vs \\
    $ \mathcal{R} $ is a simply connected region, $C$ is a simple-closed contour in $ \mathcal{R} $, with 0 the centre, and hence inside, of both $C$ and $\Gamma$. In each case our integrand is analytic on $ \mathcal{R} - \{ 0 \} $.
    \vs \\
    Thus, by the SHRINKING CONTOUR Theorem, and, by the REVERSE CONTOUR Theorem (to account for the direction), our answers would all be multiplied by -1.
\end{enumerate}
